\clearpage
\SourceLeaf{393}%
\begin{center}\bfseries [LATERAL CURVATURE.]\end{center}
\begin{center}\bfseries [1.]\end{center}
\noindent\hspace*{1.2em}The parameters [of a point $x,y,z$] on the surface [are] $t,\theta$ [and the differentials of the coordinates $x,y,z$]:
\begin{equation}\tag*{[1]}
\left\{\begin{aligned}
dx&=a\,dt+\alpha\,d\theta\\
dy&=b\,dt+\beta\,d\theta\\
dz&=c\,dt+\gamma\,d\theta
\end{aligned}\right.
\end{equation}
\begin{equation}\tag*{[2]}
dx^2+dy^2+dz^2=P\,dt^2+2Q\,dt\,d\theta+R\,d\theta^2.
\end{equation}
[If, further, we denote by] $X,Y,Z$ [the] direction cosines of the normal, then]
\begin{equation}\tag*{[3]}
\left\{\begin{aligned}
\sigma X&=b\gamma-c\beta\\
\sigma Y&=c\alpha-a\gamma\\
\sigma Z&=a\beta-b\alpha.
\end{aligned}\right.
\end{equation}

[Now consider a curve on the surface.] Let [then] $t$ and $\theta$ be functions of $u$. [The direction cosines
\[
\begin{array}{cl}
\xi,\ \eta,\ \zeta&\text{of the tangent}\\
\xi',\ \eta',\ \zeta'&\text{of the principal normal}\\
\xi'',\ \eta'',\ \zeta''&\text{of the binormal}
\end{array}
\]
then satisfy the following equations. First, we have:]
